\begin{table}[htbp]
\raggedright
\caption{Profile moderators of attacked effectivity: controlled contrasts, SEM paths, and descriptive susceptibility weights.}
\label{tab:multivariate}
{
\footnotesize
\setlength{\tabcolsep}{4pt}
\renewcommand{\arraystretch}{1.08}
\resizebox{\linewidth}{!}{%
\begin{tabular}{p{0.34\linewidth}rrrrrrrrr}
\toprule
Moderator & Role & Controlled mean b & Controlled p & Ridge mean b & Weight \% & SEM mean b & SEM mean |b| & SEM min p & SEM paths \\
\midrule
Big Five Conscientiousness Mean \% & core & -2.511 & 0.009 &  &  & -2.470 & 3.179 & 0.006 & 4.000 \\
Big Five Neuroticism Mean \% & core & -0.026 & 0.979 &  &  & 0.473 & 0.682 & 0.374 & 4.000 \\
Sex Other & core & 0.622 & 0.777 & 1.081 & 5.221 & 1.246 & 4.942 & 0.075 & 4.000 \\
Sex Female & core & 0.413 & 0.833 & -0.696 & 2.809 & 1.400 & 4.708 & 0.095 & 4.000 \\
Big Five Openness To Experience Mean \% & core & 1.395 & 0.096 &  &  & 0.899 & 1.478 & 0.179 & 4.000 \\
Big Five Extraversion Mean \% & exploratory & -0.986 & 0.340 &  &  &  &  &  &  \\
Big Five Agreeableness Mean \% & exploratory & -0.277 & 0.735 &  &  &  &  &  &  \\
\bottomrule
\end{tabular}
}
}
\par\smallskip
{\normalsize Note. Controlled coefficients summarize moderator associations with mean attacked effectivity. SEM columns summarize the repeated-outcome path model across the attacked opinion leaves. Weight percentages come from the target-conditional regularized aggregation used to compute the post hoc susceptibility index.}
\end{table}